\section{第\ref{sec:dendrites}章　樹状突起の数}

クラスの構造と入力の数は、解剖学的方法と電気的記録で測定されている。見積もり~\eqref{eq:dendriteload}はコンデンサの単純な物理である。入力数の計算上の説明は、理論とモデルに依拠している。

{\footnotesize
\setlength{\tabcolsep}{3pt}\renewcommand{\arraystretch}{1.15}
\begin{longtable}{>{\raggedright}p{0.5cm}>{\raggedright}p{4.0cm}>{\raggedright}p{5.6cm}>{\raggedright}p{2.3cm}>{\raggedright\arraybackslash}p{1.7cm}}
\toprule
№ & 主張 & 実験と測定内容 & 出典 & 状態 \\
\midrule\endfirsthead
\toprule
№ & 主張 & 実験と測定内容 & 出典 & 状態 \\
\midrule\endhead
1 & 膜の比容量は$c_m\approx1$~\textmu F/cm$^2$ & ニューロンの細胞体から膜をピペットに取り、電圧ステップを与えて容量電流の減衰から全容量を求め、面積で割った：3クラスのニューロンで$0.9$~\textmu F/cm$^2$ & \cite{Gentet2000} & 測定済み \\
2 & 充電時間$t\approx c_mA\Delta V/I$~\eqref{eq:dendriteload}：大きなニューロンには数十個の同時入力が必要で、小さなニューロンには大きなシナプス1個で足りる & コンデンサの充電法則による見積もり。数値（$20$と$300$~pF、$0.1$と数nA）は桁の値 & 本章で導出 & 理論 \\
3 & 1個の巨大シナプスは数nAの電流を流し、小さなニューロンをすぐに閾値に届かせる & スライスで終末と受け手を同時記録：シナプス1個の電流は$2$--$13$~nA、入力スパイクのたびに閾値を超える応答が生じ、遅れは$36^\circ$Cで約$0.4$~ms。受け手は\nt{GLYC}を出す & \cite{Borst1995,BorstSoria2012} & 測定済み \\
4 & 樹状突起ゼロ：触覚と痛みの感覚ニューロンでは、スパイクはセンサから細胞体を経ずに脊髄へ向かう & 構造（細胞体の近くで二つに分かれる軸索）は解剖学的に確立。伝導に細胞体の興奮性が要らないことは、このニューロンの検証済みモデルで示された：興奮性のない細胞体でも、スパイクは分岐点を同じく確実に通過する & \cite{AmirDevor2003} & 理論 \\
5 & 樹状突起1本：嗅覚ニューロンは1種類の受容体をもち、1本の入力線を伝える & 受容体遺伝子の実験の総説：各嗅覚ニューロンは大きな遺伝子ファミリーのうち一つの受容体遺伝子を発現する & \cite{Mombaerts2004} & 測定済み \\
6 & 樹状突起2本：音の方向検出器では、左耳と右耳からの入力が別々の樹状突起に接続する & 総説にまとめられた哺乳類の解剖と記録：両耳からの入力は反対向きの2本の樹状突起に分かれている & \cite{Grothe2010} & 測定済み \\
7 & 入力を2本の樹状突起に分けるとANDが厳密になる & モデルで示された：1本の樹状突起での飽和~\eqref{eq:shunt}により、片側からの入力は閾値に届かず、同時検出が改善する。入力の配置を変えた直接の実験はない & \cite{AgmonSnir1998} & 理論 \\
8 & 運動学習回路の小さな\nt{GLUT}ニューロンは約$7\cdot10^{10}$個で、それぞれ$\sim4$個の入力 & 等方性分画法による細胞計数：ヒトの脳のこの部分に約$690$億個のニューロン。生きたラットでの単一細胞記録：接触に対して少数の入力から離散的な電流のバーストが届き、それらが加算されると細胞が発火する & \cite{Azevedo2009,Chadderton2004} & 測定済み \\
9 & $n$入力の閾値素子が分けられるランダムなパターンは$P_{\max}\approx2n$個まで~\eqref{eq:cover}。運動学習回路の出力ニューロンでは上限$\sim2\cdot10^5$、現実的な見積もりは$\sim4\cdot10^4$個の対応 & 上限は組合せ幾何学の古典的な結果。現実的な見積もりは、正の重みのみとノイズへの余裕を仮定した容量理論を、このニューロンの測定された入力数に適用したもの & \cite{Cover1965,Brunel2004} & 理論 \\
10 & 入力の少ないミキサは入力の拡張に必要で、「4」は最適に近い & 理論：ミキサ層が与える表現の次元は、解剖学的に観察される入力数付近で最大になる。拡張の元の仮説は古典的な研究にある & \cite{Marr1969,Albus1971,LitwinKumar2017} & 仮説 \\
11 & 大きな木：$10^4$--$10^5$個の入力 & 電子顕微鏡像によるシナプスの計数：ラットの運動学習回路の出力\nt{GABA}ニューロン1個に$\approx1.7\cdot10^5$個のシナプス。海馬の\nt{GLUT}ニューロンに約$30\,000$個の興奮性入力と$1\,700$個の抑制性入力 & \cite{NapperHarvey1988,Megias2001} & 測定済み \\
12 & 大きな木の枝は、それぞれ閾値をもつ別々のサブ回路であり、ニューロンは小さな多層ネットワークである & 細い樹状突起の2か所を局所刺激：同じ枝の入力はシグモイド状に、異なる枝の入力は線形に加算される。ヒト皮質ニューロンの樹状突起で段階的なカルシウムスパイクが見つかり、1個の細胞が線形分離不可能な課題を解ける。モデル：1個のニューロンを再現するには深いネットワークが必要 & \cite{Polsky2004,Gidon2020,Beniaguev2021} & 測定済み \\
13 & 出力を兼ねる樹状突起：網膜の\nt{GABA}ニューロンの各枝は独立した方向検出器 & 個々の枝での二光子カルシウム記録：枝の信号は細胞体から先端への動きに選択的だが、細胞体の電圧はそうではない & \cite{Euler2002} & 測定済み \\
14 & 入力数は課題に従う（命題~\ref{prop:dendrites}） & クラスの構造は測定済み（3--13行目）。入力数が速さや分類のために選ばれたことは、個々のクラスについて理論とモデルで示されているだけである & \cite{LitwinKumar2017,AgmonSnir1998} & 仮説 \\
\bottomrule
\end{longtable}}
